% Source-by-source decisions and original pages: authoring/notes/unit02-exercise-sources.md.
% All matrices and scalars are real. Each selected source part is solved below.
\subsection*{Inverses and matrix equations}

% Bapat 2.15, both parts; no supplied Chapter 12 answer.
\begin{problem}[Two inverse calculations]\label{ex:02:bapat-2-15}
Find the inverses of
\[
\text{(a)}\quad\begin{pmatrix}a&b\\c&d\end{pmatrix},\quad ad-bc\ne0;
\qquad
\text{(b)}\quad\begin{pmatrix}2&-1&0\\2&1&-1\\1&0&4\end{pmatrix}.
\]
\end{problem}
\begin{solution}
(a) Put $\delta=ad-bc$. Multiplication in either order gives
\[
\begin{pmatrix}a&b\\c&d\end{pmatrix}
\begin{pmatrix}d&-b\\-c&a\end{pmatrix}
=\begin{pmatrix}d&-b\\-c&a\end{pmatrix}
\begin{pmatrix}a&b\\c&d\end{pmatrix}
=\delta I_2.
\]
Thus the inverse is $\delta^{-1}\begin{pmatrix}d&-b\\-c&a\end{pmatrix}$.
Division by $\delta$ is valid because $\delta\ne0$.

(b) Solve $Ax=y$ for an arbitrary $y=(y_1,y_2,y_3)^\top$. The first two equations give
\[
x_2=2x_1-y_1,\qquad x_3=2x_1+x_2-y_2=4x_1-y_1-y_2.
\]
The third becomes $17x_1=y_3+4y_1+4y_2$. Substitution yields
\[
x=\frac1{17}
\begin{pmatrix}4&4&1\\-9&8&2\\-1&-1&4\end{pmatrix}y.
\]
Writing the displayed integer matrix as $K$, direct multiplication gives
$AK=KA=17I_3$. Therefore $A^{-1}=K/17$.
\end{solution}

% MA 4.22(a,b,e); source (c,d) already proved in the exposition.
\begin{problem}[Inverse identities]\label{ex:02:inverse-identities}
Let $A$ be invertible. Prove the following statements.
\begin{enumerate}
\item If $\lambda\ne0$, then $(\lambda A)^{-1}=\lambda^{-1}A^{-1}$.
\item $(A^{-1})^{-1}=A$.
\item If $AB=BA$ and $B$ is invertible, then $AB^{-1}=B^{-1}A$.
\end{enumerate}
\end{problem}
\begin{solution}
(a) Both $(\lambda A)(\lambda^{-1}A^{-1})$ and
$(\lambda^{-1}A^{-1})(\lambda A)$ equal $I$.
\par (b) The equations $AA^{-1}=A^{-1}A=I$ also say that $A$ is the inverse of $A^{-1}$.
\par (c) Multiply $AB=BA$ on the left and on the right by $B^{-1}$:
\[
B^{-1}ABB^{-1}=B^{-1}BAB^{-1},\qquad B^{-1}A=AB^{-1}.
\]
\end{solution}

% Bapat 2.13; no supplied Chapter 12 answer.
\begin{problem}[One right-hand side is enough]\label{ex:02:bapat-2-13}
Let $A$ be $n\times n$ and fix $b\in\R^n$. Prove that $A$ is invertible
if and only if $Ax=b$ has exactly one solution.
\end{problem}
\begin{solution}
If $A$ is invertible, $A^{-1}b$ is a solution and multiplication by $A^{-1}$
shows that every solution equals it.
Conversely, let $x_0$ be the unique solution. If $Az=0$, then
$A(x_0+z)=b$, so uniqueness gives $z=0$.
Thus the homogeneous system has only the zero solution, and the square inverse
criterion (Theorem~\ref{thm:02:inverse}) gives the inverse. Concretely, elimination can have no free variable;
all $n$ columns have pivots, and scaling and upward elimination reduce $A$ to $I_n$.
\end{solution}

% Harville solutions 8.1(a,b); replace rank argument by proved excess-unknowns result.
\begin{problem}[Rectangular left and right inverses]\label{ex:02:rectangular-inverses}
Let $A$ be $m\times n$.
\begin{enumerate}
\item If $AR=I_m$ for some $R$, prove that $n\geqslant m$.
\item If $LA=I_n$ for some $L$, prove that $m\geqslant n$.
\end{enumerate}
\end{problem}
\begin{solution}
(a) If $m>n$, the system $A^\top y=0$ has $n$ equations in $m$ unknowns,
so Unit 1, Proposition~\ref{one-prop:01:excess}, gives a nonzero solution $y$.
Transposing gives $y^\top A=0$, hence
\[
y^\top=y^\top I_m=y^\top AR=0,
\]
a contradiction.
(b) If $n>m$, elimination gives a nonzero $z$ with $Az=0$.
Then $z=I_nz=LAz=0$, again a contradiction.
\end{solution}

% Harville solutions 8.2(a,b), all parts.
\begin{problem}[Matrices that are their own inverses]\label{ex:02:involutory}
A square matrix $A$ is called \emph{involutory} if $A^2=I$.
\begin{enumerate}
\item Prove that $A$ is involutory if and only if $(I-A)(I+A)=0$.
\item Prove that $A=\begin{pmatrix}a&b\\c&d\end{pmatrix}$ is involutory
if and only if either
\[
d=-a,\quad a^2+bc=1,
\qquad\text{or}\qquad
b=c=0,\quad d=a=1\text{ or }d=a=-1.
\]
\end{enumerate}
\end{problem}
\begin{solution}
(a) Distributivity gives $(I-A)(I+A)=I-A^2$.
(b) The equation $A^2=I_2$ is precisely
\[
a^2+bc=1,\qquad b(a+d)=0,\qquad c(a+d)=0,\qquad d^2+bc=1.
\]
If $a+d=0$, these reduce to $d=-a$ and $a^2+bc=1$.
If $a+d\ne0$, the middle equations force $b=c=0$. Then
$a^2=d^2=1$; the condition $a+d\ne0$ forces $a=d=1$ or $a=d=-1$.
Conversely, either displayed set of conditions satisfies all four equations,
so both families are involutory.
\end{solution}

% MA 4.28(a,b,c), all three formulas, using the original general-product proof.
\begin{problem}[Changing an invertible matrix by an outer product]\label{ex:02:outer-inverse}
Let $A$ be invertible and let $a,b$ be columns of the appropriate size.
Prove the following formulas under their stated scalar conditions:
\begin{enumerate}
\item If $a^\top A^{-1}a\ne-1$, then
\[
(A+aa^\top)^{-1}=A^{-1}
-\frac{A^{-1}aa^\top A^{-1}}{1+a^\top A^{-1}a}.
\]
\item If $a^\top A^{-1}a\ne1$, then
\[
(A-aa^\top)^{-1}=A^{-1}
+\frac{A^{-1}aa^\top A^{-1}}{1-a^\top A^{-1}a}.
\]
\item If $b^\top A^{-1}a\ne-1$, then
\[
(A+ab^\top)^{-1}=A^{-1}
-\frac{A^{-1}ab^\top A^{-1}}{1+b^\top A^{-1}a}.
\]
\end{enumerate}
\end{problem}
\begin{solution}
Prove (c) first. Put $t=b^\top A^{-1}a$ and
$X=A^{-1}-A^{-1}ab^\top A^{-1}/(1+t)$.
The denominator is nonzero. Associativity, with $t$ a scalar, gives
\begin{align*}
(A+ab^\top)X
 &=I+ab^\top A^{-1}
   -\frac{ab^\top A^{-1}+atb^\top A^{-1}}{1+t}=I,\\
X(A+ab^\top)
 &=I+A^{-1}ab^\top
   -\frac{A^{-1}ab^\top+A^{-1}atb^\top}{1+t}=I.
\end{align*}
This proves (c). Setting $b=a$ proves (a).
Setting $b=-a$ gives $t=-a^\top A^{-1}a$ and changes the numerator's sign,
which gives exactly (b).
\end{solution}

\subsection*{Blocks and column operations}
% MA 2.16, all original numerical data and requested operations.
\begin{problem}[Compatible block calculations]\label{ex:02:block-calculation}
Use the indicated partitions to find $A+B$, $AC$, and $A^\top$, where
\[
A=\left(\begin{array}{rrr|rr}1&3&-2&1&2\\6&8&0&-1&6\\\hline0&0&1&4&1\end{array}\right),
\quad
B=\left(\begin{array}{rrr|rr}1&-3&-2&4&1\\6&2&6&2&0\\\hline1&0&2&0&1\end{array}\right),
\]
\[
C=\left(\begin{array}{rr|rr}1&0&5&1\\0&2&0&0\\-1&0&3&1\\\hline3&5&0&2\\2&-1&3&1\end{array}\right).
\]
\end{problem}
\begin{solution}
Addition is entrywise within corresponding blocks:
\[
A+B=\left(\begin{array}{rrr|rr}2&0&-4&5&3\\12&10&6&1&6\\\hline1&0&3&4&2\end{array}\right).
\]
Writing the blocks as $A_{ij},C_{ij}$, the four blocks of $AC$ are
\begin{align*}
A_{11}C_{11}+A_{12}C_{21}
 &=\begin{pmatrix}3&6\\6&16\end{pmatrix}
   +\begin{pmatrix}7&3\\9&-11\end{pmatrix}
   =\begin{pmatrix}10&9\\15&5\end{pmatrix},\\
A_{11}C_{12}+A_{12}C_{22}
 &=\begin{pmatrix}-1&-1\\30&6\end{pmatrix}
   +\begin{pmatrix}6&4\\18&4\end{pmatrix}
   =\begin{pmatrix}5&3\\48&10\end{pmatrix},\\
A_{21}C_{11}+A_{22}C_{21}&=\begin{pmatrix}-1&0\end{pmatrix}
 +\begin{pmatrix}14&19\end{pmatrix}=\begin{pmatrix}13&19\end{pmatrix},\\
A_{21}C_{12}+A_{22}C_{22}&=\begin{pmatrix}3&1\end{pmatrix}
 +\begin{pmatrix}3&9\end{pmatrix}=\begin{pmatrix}6&10\end{pmatrix}.
\end{align*}
Consequently,
\[
AC=\left(\begin{array}{rr|rr}10&9&5&3\\15&5&48&10\\\hline13&19&6&10\end{array}\right),
\qquad
A^\top=\left(\begin{array}{rr|r}1&6&0\\3&8&0\\-2&0&1\\\hline1&-1&4\\2&6&1\end{array}\right).
\]
Transposition interchanges the block positions and transposes each block.
\end{solution}

% Harville solutions 8.4(a); (b) deliberately not selected.
\begin{problem}[Reading an inverse in blocks]\label{ex:02:split-inverse}
Suppose $A=[\,A_1\ A_2\,]$ is invertible, where $A_1$ has $p$ columns
and $A_2$ has $q$ columns. Partition
\[
A^{-1}=\begin{pmatrix}B_1\\B_2\end{pmatrix},
\qquad B_1:p\times(p+q),\quad B_2:q\times(p+q).
\]
Prove that
\[
B_1A_1=I_p,\quad B_1A_2=0,\quad B_2A_1=0,\quad B_2A_2=I_q,
\]
and that $A_1B_1=I_{p+q}-A_2B_2$ and $A_2B_2=I_{p+q}-A_1B_1$.
\end{problem}
\begin{solution}
The two defining inverse identities give
\[
A^{-1}A=
\begin{pmatrix}B_1A_1&B_1A_2\\B_2A_1&B_2A_2\end{pmatrix}
=\begin{pmatrix}I_p&0\\0&I_q\end{pmatrix},
\qquad
AA^{-1}=A_1B_1+A_2B_2=I_{p+q}.
\]
Comparing the four blocks proves the first four equations;
rearranging the last equation proves each of the last two.
\end{solution}

% MA 5.13(a,b), lower-block-triangular formulas only; upper formulas are in the exposition.
\begin{problem}[A lower triangular block inverse]\label{ex:02:triangular-block-inverse}
Let $C$ be $n\times m$.
\begin{enumerate}
\item Find the inverse of $\begin{pmatrix}I_m&0\\C&I_n\end{pmatrix}$.
\item More generally, if $A$ and $D$ are invertible of orders $m$ and $n$,
find the inverse of $L=\begin{pmatrix}A&0\\C&D\end{pmatrix}$.
\end{enumerate}
\end{problem}
\begin{solution}
For (b), try $X=\begin{pmatrix}A^{-1}&0\\Q&D^{-1}\end{pmatrix}$.
The lower-left block of $LX$ is $CA^{-1}+DQ$, so this block vanishes
exactly when $Q=-D^{-1}CA^{-1}$. For this choice,
\[
LX=\begin{pmatrix}I_m&0\\CA^{-1}-DD^{-1}CA^{-1}&I_n\end{pmatrix}=I_{m+n},
\]
\[
XL=\begin{pmatrix}I_m&0\\-D^{-1}CA^{-1}A+D^{-1}C&I_n\end{pmatrix}=I_{m+n}.
\]
Thus $L^{-1}=\begin{pmatrix}A^{-1}&0\\-D^{-1}CA^{-1}&D^{-1}\end{pmatrix}$.
Taking $A=I_m,D=I_n$ proves (a): replace $C$ by $-C$.
\end{solution}

% MA 6.10 and 6.11, both complete.
\begin{problem}[Column operations as right multiplication]\label{ex:02:column-operations}
\begin{enumerate}
\item Prove that an elementary column operation on an $m\times n$ matrix $A$
can be performed by right multiplication by an $n\times n$ elementary matrix.
\item For $A=\begin{pmatrix}1&2\\3&4\\5&6\end{pmatrix}$,
give the right multiplier and the resulting matrix for each separate operation:
swap the two columns; multiply column 1 by 3; add 5 times column 1 to column 2.
\end{enumerate}
\end{problem}
\begin{solution}
(a) Write $A=[\,a_1\ \cdots\ a_n\,]$.
The $j$th column of $AF$ is $\sum_k f_{kj}a_k$.
Swapping columns $i,j$ of $I_n$ therefore swaps columns $i,j$ of $A$;
the resulting $F$ is also obtained by swapping those rows of $I_n$.
Scaling column $j$ of $I_n$ by a nonzero $c$ gives a diagonal elementary
matrix and scales $a_j$ by $c$.
Finally, adding $c$ times column $i$ to column $j$ of $I_n$ inserts $c$
in position $(i,j)$. This same matrix is obtained by adding $c$ times row $j$
to row $i$ of $I_n$, and its right action replaces $a_j$ by $a_j+ca_i$.
Thus each multiplier is elementary.

(b) In the order requested,
\[
\begin{array}{c@{\qquad}c}
F&AF\\[2pt]
\begin{pmatrix}0&1\\1&0\end{pmatrix}
&\begin{pmatrix}2&1\\4&3\\6&5\end{pmatrix}\\[9pt]
\begin{pmatrix}3&0\\0&1\end{pmatrix}
&\begin{pmatrix}3&2\\9&4\\15&6\end{pmatrix}\\[9pt]
\begin{pmatrix}1&5\\0&1\end{pmatrix}
&\begin{pmatrix}1&7\\3&19\\5&31\end{pmatrix}.
\end{array}
\]
\end{solution}

\needspace{12\baselineskip}
\subsection*{Trace and matrix identities}
% MA 2.26(a-d), all parts. Harville 5.2 supplies next problem's square counterexample.
\begin{problem}[What cyclic trace permits]\label{ex:02:cyclic-trace}
\begin{enumerate}
\item For $m\times n$ matrices $A,B$, prove that
$\tr(A^\top B)=\tr(BA^\top)=\tr(AB^\top)=\tr(B^\top A)$.
\item Prove that $\tr(Aaa^\top)=a^\top Aa$ for square $A$ and a compatible column $a$.
\item Prove that $\tr(ABC)=\tr(CAB)=\tr(BCA)$, specifying the compatible sizes.
\item Must $\tr(ABC)=\tr(ACB)$ also hold?
\end{enumerate}
\end{problem}
\begin{solution}
(a) Each of the four expressions equals $\sum_{i=1}^m\sum_{j=1}^n a_{ij}b_{ij}$:
for example,
\[
\tr(A^\top B)=\sum_j\sum_i a_{ij}b_{ij},
\qquad \tr(AB^\top)=\sum_i\sum_j a_{ij}b_{ij}.
\]
The other two reverse the order of the scalar factors.
(b) Applying the rectangular cyclic identity to $(Aa)a^\top$ gives
$\tr((Aa)a^\top)=\tr(a^\top Aa)=a^\top Aa$, since the latter matrix is $1\times1$.
(c) Take $A:m\times n$, $B:n\times p$, and $C:p\times m$.
Applying the two-factor cyclic identity to $(AB)C$ and $(BC)A$ gives
\[
\tr(ABC)=\tr(CAB),\qquad \tr(BCA)=\tr(ABC).
\]
(d) No. For example, when $A:1\times2$, $B:2\times3$, and $C:3\times1$,
$ABC$ is square, whereas $ACB$ is not even defined: $AC$ has incompatible sizes.
\end{solution}

% Harville solutions 5.2(a,b), original counterexample retained.
\begin{problem}[Interchanging symmetric factors]\label{ex:02:symmetric-trace}
\begin{enumerate}
\item If $A,B,C$ are symmetric matrices of the same order, prove
$\tr(ABC)=\tr(BAC)$.
\item Show that this equality can fail without symmetry.
\end{enumerate}
\end{problem}
\begin{solution}
(a) Transposing the product and then moving its first factor to the end gives
\[
\tr(ABC)=\tr((ABC)^\top)=\tr(CBA)=\tr(BAC).
\]
(b) Take
\[
A=\begin{pmatrix}1&1\\0&0\end{pmatrix},\quad
B=\begin{pmatrix}1&0\\-1&0\end{pmatrix},\quad
C=\begin{pmatrix}1&0\\0&-1\end{pmatrix}.
\]
Here $AB=0$, so $\tr(ABC)=0$, whereas
$BAC=\begin{pmatrix}1&-1\\-1&1\end{pmatrix}$ has trace 2.
For any larger order one may append zero rows and columns; order 1 cannot give
a counterexample because real scalar multiplication commutes.
\end{solution}

% Bapat 1.7, no supplied answer. Real trace-square identity gives complete proof.
\begin{problem}[Testing a matrix by traces]\label{ex:02:bapat-1-7}
Let $A$ be $n\times n$. If $\tr(AB)=0$ for every $n\times n$ matrix $B$,
must $A$ be zero?
\end{problem}
\begin{solution}
Yes. The hypothesis permits $B=A^\top$, so
\[
0=\tr(AA^\top)=\sum_{i=1}^n\sum_{j=1}^n a_{ij}^2.
\]
Each term is nonnegative, and a positive term would make the sum positive.
Thus every $a_{ij}=0$ and $A=0$.
\end{solution}

% Bapat 1.11, all five conditions; Chapter 12 hint completed.
\begin{problem}[Five tests for symmetry]\label{ex:02:bapat-1-11}
For a square real matrix $A$, prove that the following conditions are equivalent:
\[
\begin{array}{ll}
\text{(i)}\ A=A^\top;&\text{(ii)}\ A^2=AA^\top;\\
\text{(iii)}\ \tr(A^2)=\tr(AA^\top);&\text{(iv)}\ A^2=A^\top A;\\
\text{(v)}\ \tr(A^2)=\tr(A^\top A).&
\end{array}
\]
\end{problem}
\begin{solution}
\begin{itemize}
\item Substitution of $A=A^\top$ and then taking traces give
\[
\text{(i)}\Longrightarrow\text{(ii)}\Longrightarrow\text{(iii)},\qquad
\text{(i)}\Longrightarrow\text{(iv)}\Longrightarrow\text{(v)}.
\]
\item The cyclic identity $\tr(AA^\top)=\tr(A^\top A)$ gives
(iii)$\Longleftrightarrow$(v).
\item To prove (iii)$\Longrightarrow$(i), put $D=A^\top-A$. Then
\begin{align*}
\tr(D^\top D)
 &=\tr\bigl(AA^\top-A^2-(A^\top)^2+A^\top A\bigr)\\
 &=2\tr(AA^\top)-2\tr(A^2)=0.
\end{align*}
We used $(A^\top)^2=(A^2)^\top$ and equality of a matrix's trace and its
transpose's trace. Since $\tr(D^\top D)=\sum_{i,j}d_{ij}^2$,
all entries of $D$ vanish. Hence $A=A^\top$.
\end{itemize}
The two chains are now linked in both directions.
\end{solution}

% MA 2.27, complete.
\begin{problem}[Diagonal and off-diagonal sums]\label{ex:02:trace-sums}
Let $\mathbf1=(1,\ldots,1)^\top\in\R^n$ and let
$D=\operatorname{diag}(a_{11},\ldots,a_{nn})$ for an $n\times n$ matrix $A$.
Prove that
\[
\mathbf1^\top A\mathbf1
=\mathbf1^\top D\mathbf1+\tr\bigl((\mathbf1\mathbf1^\top-I_n)A\bigr).
\]
\end{problem}
\begin{solution}
The cyclic identity and trace linearity give
\[
\tr\bigl((\mathbf1\mathbf1^\top-I_n)A\bigr)
=\tr(\mathbf1^\top A\mathbf1)-\tr(A)
=\mathbf1^\top A\mathbf1-\sum_i a_{ii}.
\]
Also $\mathbf1^\top D\mathbf1=\sum_i a_{ii}$. Substitution proves the formula:
the sum of all entries splits into diagonal and off-diagonal sums.
\end{solution}

\subsection*{Subspaces and representations}
% Bapat 1.8(iii); (i,ii) are already worked in the exposition.
\begin{problem}[A condition on four matrix entries]\label{ex:02:bapat-1-8}
Is the set of all $3\times3$ real matrices $A$ satisfying
$a_{11}+a_{13}=a_{22}+a_{31}$ a subspace of $\R^{3\times3}$?
\end{problem}
\begin{solution}
Yes. Put $g(A)=a_{11}+a_{13}-a_{22}-a_{31}$.
Then $g(0)=0$ and, by entrywise addition and scaling,
\[
g(\alpha A+\beta B)=\alpha g(A)+\beta g(B).
\]
Consequently $g(A)=g(B)=0$ implies $g(\alpha A+\beta B)=0$, which verifies
both nonemptiness and closure under linear combinations.
\end{solution}

% Harville 4.1(a,b) and MA 3.11, complete; matrix families grouped.
\begin{problem}[Matrix subspaces]\label{ex:02:matrix-subspaces}
A matrix is upper triangular when \(a_{ij}=0\) for \(i>j\), and
lower triangular when \(a_{ij}=0\) for \(i<j\).
Determine whether each set is a subspace of the corresponding real matrix space:
\begin{enumerate}
\item All upper triangular $n\times n$ matrices.
\item All nonsymmetric $n\times n$ matrices.
\item All lower triangular $3\times3$ matrices.
\item All symmetric $3\times3$ matrices.
\end{enumerate}
\end{problem}
\begin{solution}
(a) Yes. The entries below the diagonal are zero, and
$(\alpha A+\beta B)_{ij}=\alpha\,0+\beta\,0=0$ for $i>j$.
The zero matrix also has this property.
(b) No. The zero matrix is symmetric and is excluded.
This also settles $n=1$, when the proposed set is empty.
(c) Yes, by the same entry calculation as (a), now for $i<j$.
(d) Yes. The zero matrix is symmetric, and for symmetric $A,B$,
\[
(\alpha A+\beta B)^\top=\alpha A^\top+\beta B^\top=\alpha A+\beta B.
\]
\end{solution}

% MA 3.9(a-d), complete; (d) proved by coefficients, not geometry.
\begin{problem}[All subspaces of \texorpdfstring{$\R^2$}{R2}]\label{ex:02:r2-subspaces}
\begin{enumerate}
\item Show that $\{(0,t)^\top:t\in\R\}$ is a subspace of $\R^2$.
\item Is $\{(1,t)^\top:t\in\R\}$ a subspace?
\item Is $\{(x,y)^\top:x\geqslant0,\ y\geqslant0\}$ a subspace?
\item Prove that every subspace of $\R^2$ is $\{0\}$, all of $\R^2$,
or $\{ta:t\in\R\}$ for some nonzero $a\in\R^2$.
\end{enumerate}
\end{problem}
\begin{solution}
(a) The zero column belongs, and
$\alpha(0,s)^\top+\beta(0,t)^\top=(0,\alpha s+\beta t)^\top$ belongs.
(b) No: it excludes zero.
(c) No: it contains $(1,1)^\top$ but not its negative.

(d) Let $W$ be a nonzero subspace, and choose $a=(a_1,a_2)^\top\ne0$ in $W$.
Closure gives $\{ta:t\in\R\}\subseteq W$.
If equality holds, we have the third listed type.
Otherwise take $b=(b_1,b_2)^\top\in W$ that is not a multiple of $a$.
We claim $\delta=a_1b_2-a_2b_1\ne0$.
If $\delta=0$ and $a_1\ne0$, then $b=(b_1/a_1)a$.
If $\delta=0$ and $a_1=0$, then $a_2\ne0$, the equation forces $b_1=0$,
and $b=(b_2/a_2)a$. Both contradict the choice of $b$.
For any $x=(x_1,x_2)^\top\in\R^2$, set
\[
\alpha=\frac{b_2x_1-b_1x_2}{\delta},\qquad
\beta=\frac{a_1x_2-a_2x_1}{\delta}.
\]
Substitution gives $\alpha a+\beta b=x$, so $x\in W$ and $W=\R^2$.
Finally, $\{0\}$ and $\R^2$ are subspaces by the subspace test.
The third type contains $0=0a$ and is closed because
$\alpha(sa)+\beta(ta)=(\alpha s+\beta t)a$.
\end{solution}

% MA 3.13(a,b), both parts.
\begin{problem}[When decomposition is unique]\label{ex:02:sum-uniqueness}
Let $U,W$ be subspaces of one real linear space.
\begin{enumerate}
\item If $U\cap W=\{0\}$, prove that every $x\in U+W$ has a unique
representation $x=u+w$ with $u\in U,w\in W$.
\item If $U\cap W\ne\{0\}$, prove that some $x\in U+W$ has more than one such representation.
\end{enumerate}
\end{problem}
\begin{solution}
(a) Existence is the definition of $U+W$.
If $u+w=\widetilde u+\widetilde w$, then
\[
u-\widetilde u=\widetilde w-w\in U\cap W=\{0\}.
\]
Hence $u=\widetilde u$ and $w=\widetilde w$.
(b) Take $z\ne0$ in $U\cap W$. The two expressions $z=z+0$ and $z=0+z$
have different $U$-components and different $W$-components. Both are admissible.
\end{solution}

% Harville 4.6, both assertions; no basis terminology.
\begin{problem}[Adding and removing generators]\label{ex:02:generators}
Let $A_1,\ldots,A_p,B_1,\ldots,B_q$ belong to a linear space $V$ of real matrices.
\begin{enumerate}
\item If $A_1,\ldots,A_p$ span $V$, prove that the larger list also spans $V$.
\item If the larger list spans $V$ and each $B_j$ is a linear combination of
$A_1,\ldots,A_p$, prove that $A_1,\ldots,A_p$ span $V$.
\end{enumerate}
\end{problem}
\begin{solution}
(a) Any $X\in V$ is $\sum_i x_iA_i$; append zero coefficients for the $B_j$.
All combinations of the larger list remain in $V$ by its closure, so its span is exactly $V$.
(b) Write $B_j=\sum_i k_{ij}A_i$. For $X\in V$, a representation using the larger list gives
\[
X=\sum_i x_iA_i+\sum_j y_jB_j
 =\sum_i\left(x_i+\sum_j y_jk_{ij}\right)A_i.
\]
Thus every $X\in V$ belongs to the span of the $A_i$; their combinations also remain in $V$.
\end{solution}

% Harville 4.7; "spans but is not a basis" replaced by its coefficient hypothesis.
\begin{problem}[A redundant representation]\label{ex:02:nonunique-coefficients}
Suppose $A_1,\ldots,A_k$ span a real matrix space $V$, and that
\[
\sum_{i=1}^k z_iA_i=0
\quad\text{for some scalars }z_1,\ldots,z_k\text{ not all zero}.
\]
Prove that every $A\in V$ has more than one representation as a linear combination
of $A_1,\ldots,A_k$.
\end{problem}
\begin{solution}
Choose $A=\sum_i x_iA_i$, using the spanning assumption. Then
\[
A=\sum_i x_iA_i+\sum_i z_iA_i=\sum_i(x_i+z_i)A_i.
\]
At least one coefficient changes because at least one $z_i$ is nonzero.
Thus these are two different coefficient lists for the same matrix $A$.
\end{solution}

% Bapat 2.16; Chapter 12 answer retained and coefficient equations completed.
\begin{problem}[A homogeneous matrix equation]\label{ex:02:bapat-2-16}
Let $A$ be $9\times7$ and $B$ be $4\times3$. Prove that there is a nonzero
$7\times4$ matrix $X$ such that $AXB=0$.
\end{problem}
\begin{solution}
The entries of $X$ are 28 scalar unknowns. For each $1\leqslant i\leqslant9$
and $1\leqslant j\leqslant3$, the entry equation is
\[
(AXB)_{ij}=\sum_{r=1}^7\sum_{s=1}^4 a_{ir}b_{sj}x_{rs}=0.
\]
These are 27 homogeneous linear equations in those 28 unknowns.
Elimination has at most 27 pivot variables, so at least one variable is free.
Set one free variable to 1, all other free variables to 0, and solve backwards
for the pivot variables. The resulting solution is nonzero and gives the required $X$.
\end{solution}

\subsection*{Further problems on products and powers}
% MA 2.32(a,b), complete classification and original factorizations.
\begin{problem}[All square-zero matrices of order two]\label{ex:02:square-zero}
\begin{enumerate}
\item Prove that a symmetric real $2\times2$ matrix satisfying $A^2=0$ must be zero.
\item Prove that all real $2\times2$ solutions of $A^2=0$ are precisely
$A=pq^\top$, where $p,q\in\R^2$ and $p^\top q=0$. Show that there are infinitely many solutions.
\end{enumerate}
\end{problem}
\begin{solution}
Write $A=\begin{pmatrix}a&b\\c&d\end{pmatrix}$. Its square is zero exactly when
\[
a^2+bc=0,\quad b(a+d)=0,\quad c(a+d)=0,\quad d^2+bc=0.
\]
If $a+d\ne0$, the middle equations give $b=c=0$, and the other two give
$a=d=0$, a contradiction. Thus the conditions reduce to
\[
d=-a,\qquad a^2+bc=0.
\]
(a) Symmetry gives $c=b$, so $a^2+b^2=0$ and all four entries vanish.

(b) If $a=0$, then $bc=0$, and the matrix has one of the forms
\[
\begin{pmatrix}0&b\\0&0\end{pmatrix}
=\begin{pmatrix}1\\0\end{pmatrix}\begin{pmatrix}0&b\end{pmatrix},
\qquad
\begin{pmatrix}0&0\\c&0\end{pmatrix}
=\begin{pmatrix}0\\1\end{pmatrix}\begin{pmatrix}c&0\end{pmatrix}.
\]
These include the zero matrix. If $a\ne0$, then $bc=-a^2\ne0$, and
\[
A=\begin{pmatrix}a&b\\-a^2/b&-a\end{pmatrix}
 =\begin{pmatrix}1\\-a/b\end{pmatrix}\begin{pmatrix}a&b\end{pmatrix}.
\]
In each factorization $A=pq^\top$ and $q^\top p=p^\top q=0$.
Conversely, $pq^\top pq^\top=p(q^\top p)q^\top=0$ whenever this scalar is zero.
Finally, the first displayed family has one distinct solution for every real $b$.
\end{solution}

% MA 2.35 and 2.36(a-d), both complete; n>=1 supplied for Fibonacci indexing.
\begin{problem}[Powers and recurrences]\label{ex:02:recurrences}
\begin{enumerate}
\item Set $F=\begin{pmatrix}1&1\\1&0\end{pmatrix}$ and define
$x_{-1}=0$, $x_0=1$, $x_n=x_{n-1}+x_{n-2}$ for $n\geqslant1$.
Prove that
\[
F^n=\begin{pmatrix}x_n&x_{n-1}\\x_{n-1}&x_{n-2}\end{pmatrix}\qquad(n\geqslant1).
\]
\item Put $A=\begin{pmatrix}1&-1\\1&0\end{pmatrix}$ and
$B=\begin{pmatrix}0&-1\\1&-1\end{pmatrix}$.
Show that $B=A^2$ and $B^2=-A$; compute $A^2,\ldots,A^6$;
and deduce that $A^6=I$ and $B^3=I$.
\item Explain the relation between this $A$ and the recurrence
$u_n=u_{n-1}-u_{n-2}$.
\end{enumerate}
\end{problem}
\begin{solution}
(a) For $n=1$, $x_1=1$, and the displayed matrix is $F$.
If the formula holds at $n$, multiplication on the right by $F$ gives
\[
F^{n+1}=
\begin{pmatrix}x_n+x_{n-1}&x_n\\x_{n-1}+x_{n-2}&x_{n-1}\end{pmatrix}
=\begin{pmatrix}x_{n+1}&x_n\\x_n&x_{n-1}\end{pmatrix}.
\]
This proves the formula by induction.
(b) Direct products give
\[
A^2=\begin{pmatrix}0&-1\\1&-1\end{pmatrix}=B,
\qquad A^3=BA=\begin{pmatrix}-1&0\\0&-1\end{pmatrix}=-I.
\]
Multiplying successively by $A$ now gives
\[
A^4=-A=\begin{pmatrix}-1&1\\-1&0\end{pmatrix},\quad
A^5=-B=\begin{pmatrix}0&1\\-1&1\end{pmatrix},\quad A^6=I.
\]
Consequently $B^2=A^4=-A$ and $B^3=A^6=I$.
(c) Given initial values $u_0,u_{-1}$, set $z_n=(u_n,u_{n-1})^\top$.
Then, for $n\geqslant1$,
\[
z_n=Az_{n-1}
\quad\Longleftrightarrow\quad
\begin{pmatrix}u_n\\u_{n-1}\end{pmatrix}
=\begin{pmatrix}u_{n-1}-u_{n-2}\\u_{n-1}\end{pmatrix}.
\]
Iterating this equality gives $z_n=A^nz_0$, and conversely these matrix powers
satisfy the recurrence by $A^{n+1}=AA^n$.
\end{solution}

% MA 9.9*, all four parts. Original star verified in the printed exercise list.
\begin{problem}[A binomial expansion without commutativity]\label{ex:02:noncommuting-binomial}
Let $A,B$ be $n\times n$ and let $p$ be a positive integer.
For $0\leqslant j\leqslant p$, let $\binom pj$ denote the number of
$j$-element subsets of $\{1,\ldots,p\}$.
\begin{enumerate}
\item Show that $(A+B)^p=A^p(I+A^{-1}B)^p$ need not hold, even when $A$ is invertible.
\item Let $W_{p,j}$ be the sum of all products of length $p$ containing $j$
factors $B$ and $p-j$ factors $A$, in all possible orders.
Count each order separately, even when different orders give the same matrix.
For example,
\[
W_{3,1}=A^2B+ABA+BA^2,
\]
\[
W_{4,2}=A^2B^2+ABAB+AB^2A+BA^2B+BABA+B^2A^2.
\]
Prove that $(A+B)^p=\sum_{j=0}^p W_{p,j}$.
\item Prove that if $AB=BA$, then
$(A+B)^p=\sum_{j=0}^p\binom pj A^{p-j}B^j$.
\item Prove that the converse in (c) holds when $p=2$, but can fail for another positive integer $p$.
\end{enumerate}
\end{problem}
\begin{solution}
(a) Take $p=2$ and
$A=\begin{pmatrix}1&0\\0&2\end{pmatrix}$,
$B=\begin{pmatrix}0&1\\0&0\end{pmatrix}$.
Then $A^{-1}B=B$, $B^2=0$, and
\[
(A+B)^2=\begin{pmatrix}1&3\\0&4\end{pmatrix},
\qquad
A^2(I+A^{-1}B)^2=A^2(I+2B)=\begin{pmatrix}1&2\\0&4\end{pmatrix}.
\]
Without invertibility, the proposed right side need not even be defined.

(b) At $p=1$ the sum is $A+B$.
Suppose the expansion holds at $p$. Multiplication on the left by $A+B$
places either $A$ or $B$ in front of each product of length $p$.
Every order of $p+1$ factors occurs exactly once: its first factor specifies
the choice, and deleting it identifies the remaining order.
More explicitly, setting $W_{p,-1}=W_{p,p+1}=0$ gives
\[
W_{p+1,j}=AW_{p,j}+BW_{p,j-1}\quad(0\leqslant j\leqslant p+1).
\]
Distributing $(A+B)\sum_{j=0}^pW_{p,j}$ and grouping by the number of $B$ factors
therefore proves the expansion at $p+1$.

(c) Under $AB=BA$, adjacent unlike factors can be interchanged, so moving every
$A$ to the left turns each summand of $W_{p,j}$ into $A^{p-j}B^j$.
There are $\binom pj$ such summands, one for each choice of the $j$ positions
occupied by $B$. Thus $W_{p,j}=\binom pj A^{p-j}B^j$, and (b) gives the formula.

(d) For $p=2$, comparison of
\[
(A+B)^2=A^2+AB+BA+B^2
\quad\text{and}\quad A^2+2AB+B^2
\]
forces $BA=AB$. For $p=3$, take
\[
A=\begin{pmatrix}0&0&1\\0&0&1\\0&0&0\end{pmatrix},\qquad
B=\begin{pmatrix}0&1&0\\0&0&1\\0&0&0\end{pmatrix}.
\]
Multiplication gives
\[
A^2=0,\quad AB=0,\quad
BA=B^2=\begin{pmatrix}0&0&1\\0&0&0\\0&0&0\end{pmatrix}\ne0,
\quad B^3=0.
\]
Also
\[
A+B=\begin{pmatrix}0&1&1\\0&0&2\\0&0&0\end{pmatrix},\quad
(A+B)^2=\begin{pmatrix}0&0&2\\0&0&0\\0&0&0\end{pmatrix},\quad (A+B)^3=0.
\]
Since $A^3=A^2B=AB^2=B^3=0$, the claimed binomial expression for $p=3$
also equals zero, although $AB\ne BA$.
\end{solution}
